\clearpage
\VerdierPhysicalPageAnchor{0121}%
\begingroup
\small
\setlength{\emergencystretch}{1.5em}%
\setlength{\abovedisplayskip}{4pt plus 2pt minus 2pt}%
\setlength{\belowdisplayskip}{4pt plus 2pt minus 2pt}%

\VerdierSemanticAnchor{ch2-1-2-11-iv-to-i-complete}
démontrer l'implication $\mathrm{iv)}\Longrightarrow\mathrm{i)}$. Soit donc
$f\colon X\longrightarrow Y$ un morphisme muni d'un scindage
$g\colon Y\longrightarrow X$. Le foncteur noyau du morphisme
$f\colon X\longrightarrow Y$ est alors égal au foncteur noyau du morphisme
$gf\colon X\longrightarrow X$ qui est un projecteur de $X$
({\itshape i.e.} $gfgf=gf$).

\VerdierSemanticAnchor{ch2-1-2-11-projector-reduction}
On est donc ramené à démontrer que, dans une catégorie triangulée, où les
produits dénombrables sont représentables, les projecteurs admettent des
noyaux. Soit donc $e\colon X\longrightarrow X$ un projecteur de $X$.
Posons $X_i=X$, $i\in\mathbb N$, soient :

\VerdierSemanticAnchor{ch2-1-2-11-phi}
\[
  \Phi\colon\prod_{i\in\mathbb N}X_i\longrightarrow
  \prod_{i\in\mathbb N}X_i, ,
\]
le morphisme défini par :

\VerdierSemanticAnchor{ch2-1-2-11-1}
\begin{equation}
  \Phi(x_0,x_1,x_2,\ldots)
  =\bigl(ex_0+(1-e)x_1,ex_1+(1-e)x_2,\ldots\bigr)\,.
  \tag{1.2.11.1}
\end{equation}
La formule (1.2.11.1) définit un morphisme en interprétant les symboles
$x_0,x_1,\ldots$ comme des morphismes d'un objet quelconque de
$\mathcal D$ dans les $X_i$. De même, soit :

\VerdierSemanticAnchor{ch2-1-2-11-psi}
\[
  \Psi\colon\prod_{i\in\mathbb N}X_i\longrightarrow
  \prod_{i\in\mathbb N}X_i, ,
\]
le morphisme défini par :

\VerdierSemanticAnchor{ch2-1-2-11-2}
\begin{equation}
  \Psi(x_0,x_1,x_2,\ldots)
  =\bigl(x_0,(1-e)x_0+ex_1,(1-e)x_1+ex_2,\ldots\bigr)\,.
  \tag{1.2.11.2}
\end{equation}

\VerdierSemanticAnchor{ch2-1-2-11-kernel-conclusion}
Il est clair que $\Phi\Psi=\operatorname{id}$ et, par suite, le morphisme
$\Phi$ est un épimorphisme. Il admet donc, d'après les implications déjà
démontrées, un noyau représentable. D'autre part, le foncteur noyau de
$\Phi$ est canoniquement isomorphe au foncteur noyau de
$e\colon X\longrightarrow X$. Ce foncteur est donc représentable et, par
suite, le projecteur $e$ admet un noyau. Ceci achève la démonstration de la
proposition (1.2.9).

\VerdierSemanticAnchor{ch2-1-2-12}
\medskip
\noindent\textbf{1.2.12.}\quad
Le corollaire (1.2.4) montre que les triangles distingués contenant un
morphisme donné sont uniques à isomorphisme non unique près.

\VerdierSemanticAnchor{ch2-1-2-12-functorial-question}
De plus, l'axiome TRIII impose que tout morphisme de morphismes se prolonge
en un morphisme de triangles. On peut donc se demander s'il est possible de
définir la structure triangulée en associant fonctoriellement à tout
morphisme de la catégorie un triangle distingué. Nous allons montrer que
dans l'affirmative la catégorie est, sous certaines hypothèses
supplémentaires, décomposable.

\VerdierSemanticAnchor{ch2-1-2-13}
\medskip
\noindent\textbf{Proposition 1.2.13.}\quad
\textit{Soient $\mathcal D$ une catégorie triangulée et $\operatorname{Tr}$
un foncteur de $\mathrm{Fl}_1(\mathcal D)$ (chap.~I, 3.3.6), catégorie des
morphismes de $\mathcal D$, dans la catégorie des triangles distingués
(sous-catégorie pleine de la catégorie des triangles de $\mathcal D$
définie par les triangles distingués) :}

\VerdierSemanticAnchor{ch2-1-2-13-tr-assignment}
\[
  \operatorname{Tr}\!\left(X\xrightarrow{f}Y\right)
  =X\xrightarrow{f}Y\xrightarrow{a(f)}\tau(f)
  \xrightarrow{b(f)}X[1], .
\]
\endgroup
